\section{Tables for the scaffold}\label{app:tables}

This appendix lists the finite data of \cref{sec:scaffold}: Ollinger's substitution \cite{Ollinger08}, in the transcription of \cite{LeanWang}, and its ruler decoration. Every entry below was derived again from the definitions in that section; the counts are obtained by listing, not by computer search. Vertical lane atoms are $\mathsf r_0,\mathsf r_2,\mathsf g_0,\mathsf g_2$ and horizontal lane atoms are $\mathsf r_1,\mathsf r_3,\mathsf g_1,\mathsf g_3$; the letter $\mathsf r$ means red and $\mathsf g$ means green.

\begin{table}[ht]
\centering
\begin{tabular}{c|cccc}
 & $\partial_{\mathrm N}$ & $\partial_{\mathrm S}$ & $\partial_{\mathrm E}$ & $\partial_{\mathrm W}$\\\hline
thin $\mathsf a$ & $\mathsf r_0$ & $\mathsf g_2$ & $\mathsf r_1$ & $\mathsf g_3$\\
thin $\mathsf b$ & $\mathsf g_0$ & $\mathsf r_2$ & $\mathsf g_1$ & $\mathsf r_3$\\
thin $\mathsf c$ & $\mathsf r_2$ & $\mathsf g_0$ & $\mathsf g_3$ & $\mathsf r_1$\\
thin $\mathsf d$ & $\mathsf g_2$ & $\mathsf r_0$ & $\mathsf r_3$ & $\mathsf g_1$\\\hline
corner $\angle\mathsf a$ & $\mathsf g_2$ & $\mathsf r_0$ & $\mathsf r_3$ & $\mathsf g_1$\\
corner $\angle\mathsf b$ & $\mathsf r_0$ & $\mathsf g_2$ & $\mathsf r_1$ & $\mathsf g_3$\\
corner $\angle\mathsf c$ & $\mathsf r_2$ & $\mathsf g_0$ & $\mathsf g_3$ & $\mathsf r_1$\\
corner $\angle\mathsf d$ & $\mathsf g_0$ & $\mathsf r_2$ & $\mathsf g_1$ & $\mathsf r_3$\\\hline
line $\langle\alpha,\gamma\rangle$ & $\alpha$ & $\alpha$ & $\gamma$ & $\gamma$
\end{tabular}
\qquad
\begin{tabular}{c|cccc}
 & $\partial_{\mathrm N}$ & $\partial_{\mathrm S}$ & $\partial_{\mathrm E}$ & $\partial_{\mathrm W}$\\\hline
black $\mathsf A$ & $0$ & $1$ & $0$ & $1$\\
black $\mathsf B$ & $1$ & $0$ & $0$ & $0$\\
black $\mathsf C$ & $1$ & $0$ & $1$ & $1$\\
black $\mathsf D$ & $1$ & $1$ & $1$ & $0$\\
black $\mathsf E$ & $0$ & $0$ & $1$ & $0$
\end{tabular}
\caption{Endpoints of the thin and thick components, and bits of the black component. A line $\langle\alpha,\gamma\rangle$ has a vertical atom $\alpha$ and a horizontal atom $\gamma$; all sixteen lines occur. The edge colour of a tile $(\theta,\kappa,\omega)$ on a side is the triple of the three entries on that side.}
\label{scf:tab-edges}
\end{table}

The left table gives the following consequences, used in \cref{scf:base}.
\begin{itemize}
\item Each corner has exactly the four endpoints of one thin type: $F(\mathsf a)=\angle\mathsf b$, $F(\mathsf b)=\angle\mathsf d$, $F(\mathsf c)=\angle\mathsf c$, $F(\mathsf d)=\angle\mathsf a$.
\item Thin types match horizontally exactly in the pairs $(\mathsf a,\mathsf c)$, $(\mathsf c,\mathsf a)$, $(\mathsf b,\mathsf d)$, $(\mathsf d,\mathsf b)$ (west tile first), and vertically exactly in the pairs $(\mathsf a,\mathsf d)$, $(\mathsf d,\mathsf a)$, $(\mathsf b,\mathsf c)$, $(\mathsf c,\mathsf b)$ (south tile first).
\item $\partial_{\mathrm S}$ of the thin types $\mathsf a,\mathsf b,\mathsf c,\mathsf d$ is $\mathsf g_2,\mathsf r_2,\mathsf g_0,\mathsf r_0$; $\partial_{\mathrm N}$ is $\mathsf r_0,\mathsf g_0,\mathsf r_2,\mathsf g_2$; $\partial_{\mathrm W}$ is $\mathsf g_3,\mathsf r_3,\mathsf r_1,\mathsf g_1$; $\partial_{\mathrm E}$ is $\mathsf r_1,\mathsf g_1,\mathsf g_3,\mathsf r_3$.
\end{itemize}
The black maps of the substitution are
\[
\begin{array}{c|ccccc}
\omega & \mathsf A&\mathsf B&\mathsf C&\mathsf D&\mathsf E\\\hline
\operatorname{se}(\omega)&\mathsf B&\mathsf B&\mathsf C&\mathsf C&\mathsf C\\
\operatorname{nw}(\omega)&\mathsf E&\mathsf D&\mathsf D&\mathsf D&\mathsf E
\end{array}
\]
They are characterized by $\partial_{\mathrm W}\operatorname{se}(\omega)=\partial_{\mathrm E}\omega$ with $\operatorname{se}(\omega)\in\{\mathsf B,\mathsf C\}$, and $\partial_{\mathrm S}\operatorname{nw}(\omega)=\partial_{\mathrm N}\omega$ with $\operatorname{nw}(\omega)\in\{\mathsf D,\mathsf E\}$. They also satisfy $\partial_{\mathrm E}\operatorname{se}(\omega)=\partial_{\mathrm E}\omega$, $\partial_{\mathrm N}\operatorname{nw}(\omega)=\partial_{\mathrm N}\omega$, $\partial_{\mathrm S}\operatorname{se}(\omega)=0$, $\partial_{\mathrm W}\operatorname{nw}(\omega)=0$, $\partial_{\mathrm N}\operatorname{se}(\omega)=1$ and $\partial_{\mathrm E}\operatorname{nw}(\omega)=1$ for every $\omega$.

\begin{table}[ht]
\centering
\begin{tabular}{c|c|c|c|r}
thin & thick & black & condition & tiles\\\hline
$\mathsf a$ & corner & $\mathsf A$ & none & $4$\\
$\mathsf a$ & line $\langle\alpha,\gamma\rangle$ & $\mathsf B$ & $\alpha\ne\mathsf r_0$ & $12$\\
$\mathsf a$ & line $\langle\alpha,\gamma\rangle$ & $\mathsf C$ & $\alpha\ne\mathsf r_2$ & $12$\\
$\mathsf a$ & line $\langle\alpha,\gamma\rangle$ & $\mathsf D$ & $\gamma\ne\mathsf r_3$ & $12$\\
$\mathsf a$ & line $\langle\alpha,\gamma\rangle$ & $\mathsf E$ & $\gamma\ne\mathsf r_1$ & $12$\\
$\mathsf b$ & corner & $\mathsf A$ & none & $4$\\
$\mathsf c$ & line $\langle\alpha,\gamma\rangle$ & $\mathsf B$ & $\alpha\ne\mathsf r_0$ & $12$\\
$\mathsf c$ & line $\langle\alpha,\gamma\rangle$ & $\mathsf C$ & $\alpha\ne\mathsf r_2$ & $12$\\
$\mathsf d$ & line $\langle\alpha,\gamma\rangle$ & $\mathsf D$ & $\gamma\ne\mathsf r_3$ & $12$\\
$\mathsf d$ & line $\langle\alpha,\gamma\rangle$ & $\mathsf E$ & $\gamma\ne\mathsf r_1$ & $12$\\\hline
\multicolumn{4}{r|}{total} & $104$
\end{tabular}
\caption{The $104$ tiles of $\mathbb O$. A line is excluded in a black class by one of four red atoms.}
\label{scf:tab-alphabet}
\end{table}

\begin{table}[ht]
\centering
\begin{tabular}{c|c|c|c}
child & thin & thick & black\\\hline
$(0,0)$ & $\mathsf a$ & $\kappa$ & $\omega$\\
$(1,0)$ & $\mathsf c$ & $\langle\partial_{\mathrm S}\theta,\ \partial_{\mathrm E}\kappa\rangle$ & $\operatorname{se}(\omega)$\\
$(0,1)$ & $\mathsf d$ & $\langle\partial_{\mathrm N}\kappa,\ \partial_{\mathrm W}\theta\rangle$ & $\operatorname{nw}(\omega)$\\
$(1,1)$ & $\mathsf b$ & $F(\theta)$ & $\mathsf A$
\end{tabular}
\qquad
\begin{tabular}{c|c|c}
child & $\chi_1$ & $\chi_2$\\\hline
$(0,0)$ & $1-b_1$ & $1-b_2$\\
$(1,0)$ & $0$ & $1-b_2$\\
$(0,1)$ & $1-b_1$ & $0$\\
$(1,1)$ & $0$ & $0$
\end{tabular}
\caption{The substitution $\varphi$ of a tile $(\theta,\kappa,\omega)$ (left), and the bits of the decorated substitution $\varphi'$ of a letter with bits $(b_1,b_2)$ (right). Child $(1,0)$ is south-east and child $(0,1)$ is north-west.}
\label{scf:tab-subst}
\end{table}

\begin{table}[ht]
\centering
\resizebox{\textwidth}{!}{\begin{tabular}{c|c|c|c|c}
$\theta$ & blacks allowed by children $(1,0)$, $(0,1)$ & formal parents & in $\mathbb O$ & excluded\\\hline
$\mathsf a$ & all & $52$ & $52$ & $0$\\
$\mathsf b$ & $\mathsf A$ & $4$ & $4$ & $0$\\
$\mathsf c$ & $\mathsf B,\mathsf C,\mathsf D$ & $36$ & $24$ & $12$: black $\mathsf D$, line with $\gamma\ne\mathsf r_3$\\
$\mathsf d$ & $\mathsf C,\mathsf D,\mathsf E$ & $36$ & $24$ & $12$: black $\mathsf C$, line with $\alpha\ne\mathsf r_2$\\\hline
total & & $128$ & $104$ & $24$
\end{tabular}}
\caption{Formal parents of an aligned legal $2\times2$ block (\cref{scf:parents}). The count of formal parents uses the $52$ pairs $(\kappa,\omega)$ allowed at child $(0,0)$.}
\label{scf:tab-parents}
\end{table}

\begin{table}[ht]
\centering
\begin{tabular}{c|c|c|c|r}
thin & base tiles & bit rule & variants each & letters\\\hline
$\mathsf a$ & $4$ corners & $b_1=b_2$ & $2$ & $8$\\
$\mathsf a$ & $8$ red crossings & $b_1\ne b_2$ & $2$ & $16$\\
$\mathsf a$ & $40$ other lines & none & $4$ & $160$\\
$\mathsf b$ & $4$ corners & $b_1=b_2=0$ & $1$ & $4$\\
$\mathsf c$ & $4$ red crossings & $b_1=0$, $b_2=1$ & $1$ & $4$\\
$\mathsf c$ & $20$ other lines & $b_1=0$ & $2$ & $40$\\
$\mathsf d$ & $4$ red crossings & $b_1=1$, $b_2=0$ & $1$ & $4$\\
$\mathsf d$ & $20$ other lines & $b_2=0$ & $2$ & $40$\\\hline
\multicolumn{4}{r|}{total} & $276$
\end{tabular}
\caption{The $276$ letters of $\mathbb A$. A red crossing is a line whose two atoms are red. Each black class $\mathsf B,\mathsf C,\mathsf D,\mathsf E$ admits exactly two of the four red crossings (\cref{scf:tab-alphabet}), which gives $8$ red crossings with thin type $\mathsf a$ and $4$ each with thin types $\mathsf c$ and $\mathsf d$.}
\label{scf:tab-letters}
\end{table}

\begin{table}[ht]
\centering
\begin{tabular}{c|c|c|c}
position & thin & thick & red crossing exactly when\\\hline
$(2,1)$ & $\mathsf d$ & $\langle\partial_{\mathrm S}\theta,\ \mathsf r_1\rangle$ & $\theta\in\{\mathsf b,\mathsf d\}$\\
$(2,3)$ & $\mathsf d$ & $\langle\partial_{\mathrm N}\theta,\ \mathsf r_3\rangle$ & $\theta\in\{\mathsf a,\mathsf c\}$\\
$(1,2)$ & $\mathsf c$ & $\langle\mathsf r_0,\ \partial_{\mathrm W}\theta\rangle$ & $\theta\in\{\mathsf b,\mathsf c\}$\\
$(3,2)$ & $\mathsf c$ & $\langle\mathsf r_2,\ \partial_{\mathrm E}\theta\rangle$ & $\theta\in\{\mathsf a,\mathsf d\}$
\end{tabular}
\caption{Four cells of the depth-two block $\varphi^2(\theta,\kappa,\omega)$, for every base tile. Position $(2,1)$ is child $(0,1)$ of child $(1,0)$; $(2,3)$ is child $(0,1)$ of child $(1,1)$; $(1,2)$ is child $(1,0)$ of child $(0,1)$; $(3,2)$ is child $(1,0)$ of child $(1,1)$.}
\label{scf:tab-witness}
\end{table}
